\chapter{\abstractname}

Early identification of individuals with \ac{MCI} who later develop \ac{AD} may support
monitoring. We evaluate longitudinal graph representation learning on resting-state
functional-connectivity data to predict conversion from \ac{MCI} to \ac{AD}. Predicting
this conversion is more clinically relevant than classifying scans across cognitively
normal, \ac{MCI}, and \ac{AD} groups, since the patients in this cohort already carry an
\ac{MCI} diagnosis and the actionable question is which of them will progress to \ac{AD}.
We train and evaluate on a pooled cohort from \ac{ADNI} and \ac{DELCODE} comprising 248
subjects, while reserving \ac{OASIS} with 60 subjects for external evaluation.

The longitudinal classifiers evaluated in the broader project first pool each visit's
connectome graph into a single representation and then model the resulting sequence
temporally. The temporal-first configuration applies a recurrent model to the sequence of
node-level connectivity features before spatial pooling, whereas the spatial-first
comparator pools each visit-level graph before temporal modelling. We compare the
temporal-first configuration with static baselines, the pre-existing spatial-first
configuration evaluated in this thesis with and without reconstruction pretraining, and a
published graph-transformer comparator.

The temporal-first configuration reaches a pooled out-of-fold \ac{AUC} of
$0.7488 \pm 0.0033$, exceeding the frozen pretrained spatial-first comparator by 0.030
\ac{AUC} with 7.6-fold fewer trainable parameters than the comparator. At the matched
two-to-three-visit window used for a direct comparison against the graph-transformer
comparator, the temporal-first configuration scores $0.7318 \pm 0.0348$ \ac{AUC} against
$0.6207 \pm 0.0338$ for the graph-transformer comparator, a margin of
$0.0999 \pm 0.0162$ \ac{AUC}. External evaluation on \ac{OASIS} yields an \ac{AUC} close
to chance due to AD diagnosis and protocol differences.
